% 正文片段：由支持中文的主文档引入；主文档加载 amsmath, amssymb。
% 设置核对：选定的保存批次 20260926T094814-5097bbbb，非重新运行结果。
\subsection{固定 MAC 网格上的制造解时间收敛}

本实验研究在固定空间离散下比较普通 IMEX--SDIRK2、IMEX--SDIRK3
及相应增量 mr-ccSAV 格式的时间误差。计算区域为单位正方形
$\Omega=(0,1)^2$，四周施加齐次无滑移边界条件，运动黏性系数
$\nu=0.01$。空间上采用 $128\times128$ 的均匀交错 MAC 网格，
所有方法使用相同离散算子、初值和终止时间 $T=2$。
SAV 格式取阻尼参数 $\gamma=1$ 和辅助变量初值 $r_0=0$；
普通格式不使用辅助变量控制。速度和阶段压力通过耦合 Stokes 系统求解，
保存配置中的线性残差及完整散度检查阈值为 $10^{-9}$。

为排除不同空间误差对时间比较的干扰，采用网格相容的半离散制造解。
先在网格顶点采样
$\phi(x,y)=\sin^2(\pi x)\sin^2(2\pi y)$，
再以离散旋度构造面速度 $W_h=\operatorname{curl}_h\phi$。
设
\begin{equation}
 a(t)=0.05\,\mathrm{e}^{-t}\bigl(1+0.25\sin(\pi t)\bigr),
 \qquad u_h^\star(t)=a(t)W_h,\qquad p_h^\star(t)=0.
\end{equation}
这里 $u_h^\star$ 表示两个面速度分量组成的离散向量，
满足离散无散约束。记 $K_h$ 为正的离散黏性算子，
$N_h$ 为代码中的对流离散，则设置离散体力
\begin{equation}
 f_h(t)=a'(t)W_h+\nu a(t)K_hW_h+a(t)^2N_h(W_h),
\end{equation}
并取初值 $u_h(0)=a(0)W_h$。该体力由与求解器相同的 MAC 算子装配，
使上述速度成为半离散动量方程的已知解；它不是独立采样连续体力的空间精度实验。

各方法分别使用固定步长
$\tau_k=0.1\,2^{-k}$，其中 $k=-1,0,\ldots,8$。
除时间格式及步长外，其余设置保持相同，总计四种格式与十个步长层。
在 $T=2$ 比较保存的数值速度与 $u_h^\star(T)$，
以 $e_h=u_h(T)-u_h^\star(T)$ 定义绝对误差
\begin{equation}
 E_{L^2}=\sqrt{(e_h,e_h)_h},\qquad
 E_{H^1}=\sqrt{(e_h,K_he_h)_h}.
\end{equation}
其中 $(\cdot,\cdot)_h$ 为 MAC 自由面速度的离散内积，
$E_{H^1}$ 是由黏性算子定义的离散半范数，而非压力误差。
对相邻步长层，在两个误差均有限且为正时记录观察阶
$p_k=\log_2(E_{k-1}/E_k)$；失败成员及缺少有效终点的成员不参与有效观察阶计算。

所选误差表的二阶格式在细步长下呈现约二阶的观察阶，两种对应方法的曲线接近。
三阶格式的 $L^2$ 观察阶在 $k=6,7$ 时约为 $2.96$ 和 $2.98$；
最细的 $k=8$ 层误差约为 $(1.20\text{--}1.22)\times10^{-13}$，
对应观察阶降至约 $1.63\text{--}1.66$，而同层 $H^1$ 半范数观察阶仍约为
$2.93\text{--}2.95$。该现象提示最细 $L^2$ 误差已接近精度受限区，
不能单凭这一层断言格式整体降阶，也不能把接近三阶的局部区间推广到全部步长。
现有材料未单独分离线性求解、时间累积及舍入误差的贡献。

结果展示直接重绘既有误差表，并保留显式保存批次、成员及报告的来源与校验信息，
不重新推进方程或重算场误差。比较必须同时保留粗步长区和细步长区：
粗步长可能尚未进入时间渐近区，细步长误差可能受线性求解容差和舍入误差限制。
因此不能要求每一层均显示形式阶数；单一空间网格的时间误差表也不能替代
空间收敛研究、无滑移 PDE 高阶收敛证明或长期稳定性验证。
